\section{人工评估与安全：平均质量之外的真实风险}

自动指标适合大量实验迭代，但多语系统的最终对象是人。不同语言评审者可能使用不同尺度，reference-based score 也无法判断所有语义、礼貌和文化问题。课堂因此把 human evaluation 与 toxicity 分成独立工作流。

\subsection{Automatic metrics 快，human evaluation 才接近真实使用}

本节进一步追问自动结果能否代表人类体验。课堂 slide 用一句对比概括两者分工：自动指标快速、适合研究迭代；人工评估昂贵、缓慢，却更接近真实质量。这里不是二选一，而是用自动指标筛选方案，再用人工判断校准最终结论。

\lecturefigure{slide-34-human-evaluation.jpg}{自动指标用于快速迭代，人工评估用于真实质量判断}{Stanford 课堂视频 00:33:45}

\begin{knowledgebox}{读图：速度与效度是两个轴}
Automatic metrics 能对四万方向快速计算，并支持 ablation；human evaluation 能判断语义保真、自然度和可接受性，却受成本、评审标准与语言覆盖限制。成熟流程需要自动指标提供 breadth，人工评估提供 validity。
\end{knowledgebox}

\begin{warningbox}{自动分数的小幅提升未必可感知}
BLEU/chrF++ 增加可能来自词形或参考匹配，并不必然改善用户体验。尤其在分数接近时，必须检查置信区间、语言分组和人工偏好，而不是把小数点后的差异写成产品级突破。
\end{warningbox}

\teachervoice{Fan 直接把 human evaluation 称为 real deal。她不是否定 automatic metrics，而是提醒：研究 iteration 的代理目标不能自动升级为真实使用结论。}

\subsection{跨语言评审尺度需要校准}

`Consistent Human Evaluation` 研究要解决的问题是：不同语言对的评审者如何给出可比较分数。若某个语言社区普遍更严格，简单平均会把审稿尺度差异误当成模型差异；因此需要共同任务说明、培训、质量检查和适当标准化。

\lecturefigure{slide-35-human-evaluation-paper.jpg}{跨语言对的一致人工机器翻译评估}{Stanford 课堂视频 00:34:27}

\begin{knowledgebox}{读图：论文标题强调 consistent，而非只强调 human}
人工评估并非天然客观。评审者需要理解源句和目标句，区分 adequacy 与 fluency，并遵守共同 rubric。跨语言一致性决定不同方向能否放在同一图中比较；否则语言间柱高差可能混入评审文化与尺度差异。
\end{knowledgebox}

设方向 $d$ 有 $n_d$ 个有效人工评分 $r_{d,i}$，最简单均值为
\[
\bar r_d=\frac{1}{n_d}\sum_{i=1}^{n_d}r_{d,i}.
\]
$r_{d,i}$ 是第 $i$ 个评分，$n_d$ 是该方向有效样本量，$\bar r_d$ 是方向平均。实际比较还需要不确定性、评审者效应与样本难度校正；只报 $\bar r_d$ 会隐藏方差。

\subsection{读人工结果：into、out-of 与 non-English}

完成评审尺度校准后，本节再读最终人工结果。课堂柱状图把三类方向放在一起。最终 progressive state 展示不同语言或语言组的人工质量分布，目的是检查自动指标中的提升是否转化为人可感知改善，以及哪些方向仍显著落后。

\lecturefigure{slide-36-human-evaluation-results.jpg}{NLLB 在三类翻译方向上的人工评估结果}{Stanford 课堂视频 00:36:03}

\begin{knowledgebox}{读图：先看组内差异，再看组间均值}
绿色、蓝色等柱对应 into-English、non-English、out-of-English 分组。先找每组内部的低谷，判断弱点是否集中在目标端低资源生成；再比较整体位置。图支持 NLLB 获得广泛人工改善，但同样显示语言间差异远未消失。
\end{knowledgebox}

\begin{warningbox}{柱状图不能证明因果}
人工结果来自完整系统，无法单独归因于 LASER3、mining、MoE、EOM 或 curriculum。要判断组件贡献必须结合 controlled ablation；最终系统图只回答联合方案是否在评测分布上有效。
\end{warningbox}

\subsection{Toxicity：不是所有翻译错误都同样严重}

本节从平均质量进一步转向尾部安全风险。质量平均分可能把灾难性错误淹没。课堂举出 COVID 场景中 `wash hands` 被译成 `hold hands` 的例子：两个短语词面接近，但行为含义相反。对公共卫生信息，这种错误的代价远高于轻微语序不自然。

\lecturefigure{slide-37-toxicity-error-severity.jpg}{公共卫生例子说明翻译错误具有不同严重度}{Stanford 课堂视频 00:37:12}

\begin{knowledgebox}{读图：错误严重度取决于语境与行动后果}
普通 metric 可能把一词差异视为局部 mismatch，但使用者会据此采取完全不同动作。安全评估需要识别 meaning inversion、added toxicity、遗漏否定和身份攻击等高代价模式，并按领域单独报告，而不是只依赖平均 BLEU。
\end{knowledgebox}

\begin{importantbox}{质量与安全必须双轴报告}
可把方向 $d$ 的验收写成二维条件：
\[
\text{Accept}(d)=\mathbf{1}[q_d\ge\tau_q]\cdot
\mathbf{1}[a_d\le\tau_a],
\]
其中 $q_d$ 是一般质量，$a_d$ 是 added-toxicity 或其他高风险错误率，$\tau_q$ 是质量下限，$\tau_a$ 是风险上限。课堂没有规定统一阈值；公式用于说明高平均质量不能抵消安全失败。
\end{importantbox}

\teachervoice{Fan 问“为什么我这么关心 toxicity”，随后用公共卫生误译说明：翻译系统的错误不是等价 token mismatch，而会改变真实行为。安全要求必须进入模型发布前评估。}

\subsection{Toxicity-200：文化相关的检测资源}

有了高风险错误例子，接下来需要建设可扩展检测资源。团队为两百语言收集 toxicity lists，并强调这些词表具有文化特异性。相同词在不同社区的攻击性、语境与禁忌可能不同，直接翻译英语脏词表会同时产生漏报和误报。

\lecturefigure{slide-38-toxicity-lists.jpg}{两百语言的文化特异 toxicity lists 用于检测与缓解}{Stanford 课堂视频 00:38:15}

\begin{knowledgebox}{读图：词表是探针，不是完整安全模型}
Toxicity list 可以检测源句没有、译文却新增的有害词，适合大规模 screening；但它无法理解引用、否定、重夺用语或上下文，也可能漏掉不在词表中的攻击表达。正确定位是 risk probe，与人工审查和语义评估配合。
\end{knowledgebox}

\begin{warningbox}{不要把英语 safety taxonomy 直接平移}
文化、身份类别和攻击方式跨语言变化。词表收集必须由母语者与社区知识校准，并记录版本与适用范围；否则“统一安全标准”可能把本地正常表达误判为有害，或漏掉真正风险。
\end{warningbox}

\subsection{本章小结}

自动指标提供迭代 breadth，人工评估提供真实质量 validity，toxicity 则检查平均分无法表达的高代价错误。三者必须并列报告，且人工尺度与文化词表都需要语言本地化校准。

